\section*{अध्याय \ref{ch:g12:seq} --- अनुक्रम}

\begin{solution}{exo:g12:seq:1}
मान लीजिए $P(n)$ कथन $\sum_{k=1}^{n} k^2 = \frac{n(n+1)(2n+1)}{6}$ है।
\emph{आधार-स्थिति:} $n = 0$ के लिए दोनों पक्ष $0$ हैं (रिक्त योग)।
\emph{आगमन-चरण:} मान लीजिए $P(n)$ है। तब
\[
\sum_{k=1}^{n+1} k^2 = \frac{n(n+1)(2n+1)}{6} + (n+1)^2
= \frac{(n+1)\bigl(n(2n+1) + 6(n+1)\bigr)}{6}
= \frac{(n+1)(2n^2 + 7n + 6)}{6}.
\]
चूँकि $2n^2 + 7n + 6 = (n+2)(2n+3)$ है, यह $\frac{(n+1)(n+2)(2(n+1)+1)}{6}$ है, अर्थात् $P(n+1)$। आगमन द्वारा $P(n)$ सभी $n$ के लिए सही
है।
\end{solution}

\begin{solution}{exo:g12:seq:2}
$a_{n+1} - a_n = \frac{n+2}{n+1} - \frac{n+1}{n}
= \frac{n(n+2) - (n+1)^2}{n(n+1)} = \frac{-1}{n(n+1)} < 0$: $(a_n)$ निरंतर \omterm{def:g10:functions:variations}{ह्रासमान} है।

$(b_n)$ के पद धनात्मक हैं और $\frac{b_{n+1}}{b_n} = \frac{2^{n+1}}{n+1}\cdot\frac{n}{2^n} =
\frac{2n}{n+1} \geq 1 \iff 2n \geq n+1 \iff n \geq 1$: $(b_n)$ \omterm{def:g12:seq:monotonic}{वर्धमान} है ($n \geq 2$ के लिए निरंतर)।

$c_{n+1} - c_n = (n+1)^2 - 10(n+1) - n^2 + 10n = 2n - 9$, जो $n \leq 4$ के लिए ऋणात्मक और $n \geq 5$ के लिए धनात्मक है: $(c_n)$ $c_5 = -25$ तक घटता है,
जो उसका \omterm{def:g10:functions:extrema}{न्यूनतम} है, और फिर बढ़ता है। वह \omterm{def:g12:seq:monotonic}{एकदिष्ट} नहीं है।
\end{solution}

\begin{solution}{exo:g12:seq:3}
हावी पद बाहर निकालिए:
\[
u_n = \frac{n^2(3 - 1/n + 1/n^2)}{n^2(2 + 5/n^2)} \xrightarrow[n\to+\infty]{} \frac{3}{2}.
\]
संयुग्मी से गुणा कीजिए:
\[
v_n = \frac{(n+1) - n}{\sqrt{n+1} + \sqrt{n}} = \frac{1}{\sqrt{n+1}+\sqrt{n}}
\xrightarrow[n\to+\infty]{} 0.
\]
अंश और हर को $3^n$ से भाग दीजिए:
\[
w_n = \frac{(2/3)^n - 1}{1 + (1/3)^n} \xrightarrow[n\to+\infty]{} \frac{0-1}{1+0} = -1,
\]
जहाँ $\abs{q} < 1$ के लिए $\lim q^n = 0$ का उपयोग हुआ है।
\end{solution}

\begin{solution}{exo:g12:seq:4}
$u_n = 5 + 3n \to +\infty$ और $v_n = 8 \cdot (1/2)^n = 2^{3-n} \to 0$। योग हैं
\[
\sum_{k=0}^{n} u_k = (n+1)\,\frac{5 + (5+3n)}{2} = \frac{(n+1)(10+3n)}{2}
\xrightarrow[n\to+\infty]{} +\infty,
\]
\[
\sum_{k=0}^{n} v_k = 8\,\frac{1 - (1/2)^{n+1}}{1 - 1/2}
= 16\left(1 - \left(\tfrac{1}{2}\right)^{n+1}\right)
\xrightarrow[n\to+\infty]{} 16 .
\]
\end{solution}

\begin{solution}{exo:g12:seq:5}
चूँकि $-1 \leq \cos n \leq 1$ है,
\[
\frac{n-1}{n+1} \leq \frac{n + \cos n}{n+1} \leq 1,
\]
और $\frac{n-1}{n+1} \to 1$, इसलिए संपीडन प्रमेय से सीमा $1$ है।

दूसरी सीमा के लिए लिखिए
\[
0 \leq \frac{n!}{n^n}
= \frac{1}{n}\cdot\frac{2}{n}\cdots\frac{n}{n}
\leq \frac{1}{n},
\]
क्योंकि $2 \leq k \leq n$ वाला हर गुणनखंड $\frac{k}{n}$ अधिक से अधिक $1$ है। चूँकि $\frac1n \to 0$ है, संपीडन
प्रमेय $\frac{n!}{n^n} \to 0$ दे देती है। (सुझाया गया \omterm{def:g12:seq:arith-geom}{गुणोत्तर} परिबंध भी चलता है: $k \leq n/2$ वाला हर
गुणनखंड अधिक से अधिक $\frac12$ है, जिससे और प्रबल परिबंध $(1/2)^{\floor{n/2}}$ मिलता है।)
\end{solution}

\begin{solution}{exo:g12:seq:6}
\emph{1.} $u_0 = 0 \in \intcc{0}{2}$। यदि $0 \leq u_n \leq 2$ हो, तो $2 \leq u_n + 2 \leq 4$, इसलिए $\sqrt{2} \leq u_{n+1} \leq 2$; विशेष रूप से $0 \leq u_{n+1} \leq 2$। आगमन से
यह गुण सभी $n$ के लिए सही है।

\emph{2.} $u_{n+1} - u_n = \sqrt{u_n + 2} - u_n$। $x \in \intcc{0}{2}$ के लिए $\sqrt{x+2} \geq x \iff x + 2 \geq x^2 \iff (2-x)(x+1) \geq 0$, जो सत्य है। अतः $(u_n)$ \omterm{def:g12:seq:monotonic}{वर्धमान} है।

\emph{3.} \omterm{def:g12:seq:monotonic}{वर्धमान} और $2$ से \omterm{def:g12:seq:bounded}{ऊपर परिबद्ध} होने के कारण $(u_n)$ किसी $\ell \in \intcc{0}{2}$ पर
\omterm{def:g12:seq:limit}{अभिसरित} होता है। $u_{n+1} = \sqrt{u_n + 2}$ में सीमा लेने पर (प्रतिचित्रण $x \mapsto \sqrt{x+2}$ संतत है) $\ell = \sqrt{\ell + 2}$ मिलता
है, इसलिए $\ell^2 - \ell - 2 = 0$, अर्थात् $\ell \in \{-1, 2\}$। चूँकि $\ell \geq 0$ है, इसलिए $\lim u_n = 2$।
\end{solution}

\begin{solution}{exo:g12:seq:7}
\emph{1.} दो ख़ुराकों के बीच दवा का $40\%$ निकल जाता है, इसलिए मात्रा $u_n$ $0.6\,u_n$
बन जाती है; और अगली ख़ुराक $1$ मात्रक जोड़ देती है: $u_{n+1} = 0.6\,u_n + 1$।

\emph{2.} $v_{n+1} = u_{n+1} - 2.5 = 0.6\,u_n + 1 - 2.5 = 0.6(u_n - 2.5)
= 0.6\,v_n$: $(v_n)$ अनुपात $0.6$ और पहले पद $v_0 = 1 - 2.5 = -1.5$ वाला \omterm{def:g11:seq:geometric}{गुणोत्तर अनुक्रम} है।
अतः $v_n = -1.5 \times 0.6^n$ और
\[
u_n = 2.5 - 1.5 \times 0.6^n .
\]

\emph{3.} चूँकि $0.6^n \to 0$ है, इसलिए $u_n \to 2.5$: दवा की मात्रा $2.5$ मात्रकों पर स्थिर हो
जाती है।
\end{solution}

\begin{solution}{exo:g12:seq:8}
\emph{1.} $u_0 = 3 > 1$। यदि $u_n > 1$ हो, तो $u_n + 2 > 0$ और
\[
u_{n+1} - 1 = \frac{4u_n - 1 - u_n - 2}{u_n + 2} = \frac{3(u_n - 1)}{u_n + 2} > 0 .
\]
आगमन से सभी $n$ के लिए $u_n > 1$ (और विशेष रूप से $u_n + 2 \neq 0$, इसलिए \omterm{def:g12:seq:sequence}{अनुक्रम} भली-भाँति
परिभाषित है)।

\emph{2.} ऊपर वाली सर्वसमिका का उपयोग करने पर
\[
v_{n+1} = \frac{1}{u_{n+1} - 1} = \frac{u_n + 2}{3(u_n - 1)}
= \frac{(u_n - 1) + 3}{3(u_n - 1)} = \frac{1}{3} + v_n .
\]
इसलिए $(v_n)$ \omterm{def:g11:seq:arithmetic}{सार्व अंतर} $\frac13$ वाला \omterm{def:g11:seq:arithmetic}{समांतर अनुक्रम} है और $v_0 = \frac{1}{u_0 - 1} = \frac12$।

\emph{3.} $v_n = \frac12 + \frac{n}{3}$, अतः $u_n = 1 + \frac{1}{v_n} = 1 + \frac{6}{3 + 2n}$। चूँकि $v_n \to +\infty$ है, इसलिए $u_n \to 1$।
\end{solution}

\begin{solution}{exo:g12:seq:9}
\emph{1.} $H_{2n} - H_n = \sum_{k=n+1}^{2n} \frac{1}{k}$ $n$ पदों का योग है, और हर पद कम से कम $\frac{1}{2n}$ है; अतः $H_{2n} - H_n \geq n \cdot \frac{1}{2n} = \frac12$।

\emph{2.} $k$ पर आगमन से: $H_{2^0} = H_1 = 1 \geq 1$। यदि $H_{2^k} \geq 1 + \frac{k}{2}$ हो, तो $n = 2^k$ के साथ बिंदु 1
लगाने पर
\[
H_{2^{k+1}} \geq H_{2^k} + \frac12 \geq 1 + \frac{k+1}{2}.
\]
\omterm{def:g12:seq:sequence}{अनुक्रम} $(H_n)$ \omterm{def:g12:seq:monotonic}{वर्धमान} है (हर क़दम $\frac{1}{n+1} > 0$ जोड़ता है) और उपअनुक्रम $H_{2^k}$
अपरिबद्ध है, इसलिए $(H_n)$ \omterm{def:g12:seq:bounded}{ऊपर परिबद्ध} नहीं है। \omterm{def:g12:seq:monotonic}{वर्धमान} और अपरिबद्ध होने के
कारण वह $+\infty$ की ओर जाता है (\cref{thm:g12:seq:monotone})।
\end{solution}

\begin{solution}{exo:g12:seq:10}
\emph{1.} \omterm{def:g12:seq:sequence}{अनुक्रम} $d_n = b_n - a_n$ $d_{n+1} - d_n = (b_{n+1} - b_n) - (a_{n+1} - a_n) \leq 0$ को संतुष्ट करता है, इसलिए $(d_n)$ \omterm{def:g10:functions:variations}{ह्रासमान} है;
और चूँकि $d_n \to 0$ है, इसलिए सभी $n$ के लिए $d_n \geq 0$ मिलता है (किसी ऋणात्मक पद
वाला \omterm{def:g10:functions:variations}{ह्रासमान} \omterm{def:g12:seq:sequence}{अनुक्रम} सदा उससे नीचे ही रहता, जिससे सीमा $0$ असंभव हो
जाती)। अतः $a_n \leq b_n$।

\emph{2.} $a_n \leq b_n \leq b_0$ से \omterm{def:g12:seq:monotonic}{वर्धमान} \omterm{def:g12:seq:sequence}{अनुक्रम} $(a_n)$ \omterm{def:g12:seq:bounded}{ऊपर परिबद्ध} है, इसलिए वह किसी $\ell$
पर \omterm{def:g12:seq:limit}{अभिसरित} होता है। इसी तरह $(b_n)$, जो \omterm{def:g10:functions:variations}{ह्रासमान} है और $a_0$ से \omterm{def:g12:seq:bounded}{नीचे परिबद्ध},
किसी $\ell'$ पर \omterm{def:g12:seq:limit}{अभिसरित} होता है। तब $\ell' - \ell = \lim (b_n - a_n) = 0$, इसलिए $\ell = \ell'$।

\emph{3.} $(a_n)$ (निरंतर) \omterm{def:g12:seq:monotonic}{वर्धमान} है, क्योंकि $a_{n+1} - a_n = \frac{1}{(n+1)!} > 0$। $(b_n)$ के लिए
\[
b_{n+1} - b_n = \frac{1}{(n+1)!} + \frac{1}{(n+1)(n+1)!} - \frac{1}{n\,n!}
= \frac{n(n+1) + n - (n+1)^2}{n(n+1)(n+1)!}
= \frac{-1}{n(n+1)(n+1)!} < 0 ,
\]
इसलिए $(b_n)$ \omterm{def:g10:functions:variations}{ह्रासमान} है। अंत में $b_n - a_n = \frac{1}{n\,n!} \to 0$। दोनों \omterm{def:g12:seq:sequence}{अनुक्रम} संलग्न हैं, अतः एक ही
सीमा पर \omterm{def:g12:seq:limit}{अभिसरित} होते हैं।
\end{solution}

\begin{solution}{pb:g12:seq:1}
\textbf{1.} $n = 1$ के लिए सत्य ($1 = 1^2$)। यदि $1 + 3 + \dots + (2n - 1) = n^2$ हो, तो अगली विषम संख्या
जोड़ने पर: $n^2 + (2n + 1) = (n + 1)^2$: वंशानुगति। आगमन से सभी $n \geq 1$ के लिए सत्य।

\textbf{2.} $2^0 = 1 > 0$। यदि $2^n > n$ हो, तो $n \geq 1$ के लिए $2^{n+1} = 2 \cdot 2^n > 2n \geq n + 1$ (और $n = 0$ सीधे जाँच
लिया जाता है): वंशानुगति, हो गया।

\textbf{3.} $n = 0$: $1 \geq 1$। यदि $(1 + x)^n \geq 1 + nx$ हो, तो $1 + x \geq 1 > 0$ से गुणा कीजिए: $(1 + x)^{n+1} \geq (1 + nx)(1 + x) = 1 + (n + 1)x + nx^2
\geq 1 + (n + 1)x$।

\textbf{4.} $n = 1$ से $n = 2$ तक का चरण: दो कंचों में ``पहले $n$'' और
``अंतिम $n$'' दो \emph{असंयुक्त} अकेले कंचे हैं --- कोई उभयनिष्ठ कंचा दोनों
समूहों को नहीं जोड़ता, इसलिए कुछ भी उनके रंगों को मिलाने पर मजबूर नहीं करता।
वंशानुगति वाला तर्क चुपके से यह माँगता है कि दोनों समूह एक-दूसरे को ढकें, जो
केवल $n \geq 2$ से आगे सही है; और आधार-स्थिति $n = 1$ के साथ शृंखला शुरू ही नहीं
हो पाती।

\textbf{5.} $4^0 - 1 = 0 = 3 \times 0$। यदि $4^n - 1 = 3k$ हो, तो $4^{n+1} - 1 = 4(4^n - 1) + 3 = 3(4k + 1)$: वंशानुगति।

\textbf{6.} $x_1 = \frac32$, $x_2 = \frac{17}{12}$, $x_3 = \frac{577}{408}$।

\textbf{7.} $x_{n+1}^2 - 2 = \frac{(x_n^2 + 2)^2 - 8x_n^2}
{4x_n^2} = \frac{(x_n^2 - 2)^2}{4 x_n^2}$: किसी धनात्मक पर एक वर्ग, अतः $\geq 0$, और जब भी $x_n^2 \neq 2$ हो तब
$> 0$। आगमन: $x_0^2 = 4 > 2$ के साथ $x_0 = 2 > 0$; यदि $x_n > 0$ और $x_n^2 > 2$ हों, तो $x_{n+1}$
(धनात्मक संख्याओं का औसत) धनात्मक है और $x_{n+1}^2 - 2 > 0$।

\textbf{8.} प्रश्न 7 से $x_{n+1} - x_n = \frac{2 - x_n^2}{2x_n} < 0$: निरंतर \omterm{def:g10:functions:variations}{ह्रासमान}।

\textbf{9.} \omterm{def:g10:functions:variations}{ह्रासमान} और \omterm{def:g12:seq:bounded}{नीचे परिबद्ध} होने के कारण ($1$ से, क्योंकि $x_n^2 > 2 > 1$
और $x_n > 0$): \omterm{def:g12:seq:monotonic}{एकदिष्ट} अभिसरण प्रमेय से $(x_n)$ किसी $L \geq 1$ पर \omterm{def:g12:seq:limit}{अभिसरित} होता है।

\textbf{10.} सीमाएँ बीजगणित का आदर करती हैं: $x_{n+1} = \frac12\left(x_n + \frac{2}{x_n}\right)$ और $x_n \to L \geq 1 > 0$ से $L = \frac12\left(L + \frac2L\right)$, इसलिए
$L^2 = 2$, और $L$ के धनात्मक होने से $L = \sqrt2$। फ़ैसला: अभिसरण सिद्ध, सीमा पहचानी
गई --- हीरोन सम्मान सहित बरी।

\textbf{11.} $x_{n+1} - \sqrt2 = \frac{x_n^2 - 2\sqrt2\,x_n +
2}{2x_n} = \frac{(x_n - \sqrt2)^2}{2x_n}$: ठीक $e_{n+1} = \frac{e_n^2}{2x_n}$, और $x_n > \sqrt2$ से $e_{n+1} \leq \frac{e_n^2}{2\sqrt2}$ मिलता है। वर्ग हो जाने वाली
त्रुटि: हर क़दम सही दशमलव स्थानों की संख्या दुगुनी कर देता है, जैसा कक्षा 9
से देखा जा रहा है।

\textbf{12.} $e_0 \approx 0.5858$, $e_1 \approx 0.0858$, $e_2 \approx 0.00245$, $e_3 \approx 2.1 \times 10^{-6}$। अनुपात $\frac{e_1}{e_0^2} \approx 0.25 = \frac{1}{2x_0}$; $\frac{e_2}{e_1^2} \approx 0.333 = \frac{1}{2x_1}$; $\frac{e_3}{e_2^2} \approx 0.353 \approx \frac{1}{2x_2}$: प्रमेय
काम पर।

\textbf{13.} $H_{2^{18}} \geq 1 + 9 = 10$: केवल $10$ पार करने के लिए ही लगभग $260\,000$ पद ($2^{18} = 262\,144$) ---
घिसट-घिसटकर अपसरण (और $H_n > 100$ के लिए किसी भी पुस्तकालय के परमाणुओं से अधिक पद
चाहिए होते)।

\textbf{14.} $S_n = 2 - \frac{1}{2^n}$ (\omterm{def:g12:seq:arith-geom}{गुणोत्तर} योग), और $\frac{1}{2^n} \to 0$ (\cref{thm:g12:seq:geometric}), इसलिए $S_n \to 2$: अनंत बार
कुतरी गई चॉकलेट पूरी की ओर जाती है पर उस तक कभी पहुँचती नहीं --- अब सीमाओं की
औपचारिक भाषा में।

\textbf{15.} $k \geq 2$ के लिए: $\frac{1}{k^2} \leq \frac{1}{k(k-1)} = \frac{1}{k-1} -
\frac1k$, इसलिए $S_n \leq 1 + \left(1 - \frac1n\right) < 2$: \omterm{def:g12:seq:monotonic}{वर्धमान} और \omterm{def:g12:seq:bounded}{ऊपर परिबद्ध}, अतः
अभिसारी (\omterm{def:g12:seq:monotonic}{एकदिष्ट} अभिसरण)। बाद में ऑयलर ने सीमा को नाम दिया: $\frac{\pi^2}{6}$।

\textbf{16.} पदों का $0$ की ओर जाना योगों के जमने के लिए \emph{आवश्यक} है
पर कुछ तय नहीं करता: हरात्मक पद $\frac1n \to 0$ फिर भी योग फट पड़ते हैं; और पद $\frac{1}{n^2} \to 0$
तथा योग \omterm{def:g12:seq:limit}{अभिसरित} हो जाते हैं। पद \emph{कितनी तेज़ी} से मरते हैं, यही पूरा
प्रश्न है --- श्रेणियों का सिद्धांत, जो स्नातक खंडों में बनता है।

\textbf{17.} $a_1 = \sqrt2 \approx 1.41421$, $b_1 = 1.5$; $a_2 \approx 1.45648$, $b_2 \approx 1.45711$: दो ही पुनरावृत्तियाँ तीन दशमलव
स्थानों तक मेल खा जाती हैं --- चकित कर देने वाली गति।

\textbf{18.} $b_{n+1} - a_{n+1} = \frac{a_n + b_n}{2} -
\sqrt{a_n b_n} = \frac{(\sqrt{b_n} - \sqrt{a_n})^2}{2} \geq
0$: \omterm{def:g11:stat:mean}{माध्य} क्रम में बने रहते हैं। $(a_n)$ बढ़ता है: $a_{n+1} = \sqrt{a_n b_n} \geq \sqrt{a_n \cdot a_n} = a_n$; और
$(b_n)$ सममित रूप से घटता है।

\textbf{19.} $\frac{b_{n+1} - a_{n+1}}{b_n - a_n} =
\frac{(\sqrt{b_n} - \sqrt{a_n})^2}
{2(\sqrt{b_n} - \sqrt{a_n})(\sqrt{b_n} + \sqrt{a_n})}
= \frac{\sqrt{b_n} - \sqrt{a_n}}{2(\sqrt{b_n} + \sqrt{a_n})}
\leq \frac12$: अंतर कम से कम आधा हो जाता है, इसलिए $b_n - a_n \to 0$; और प्रश्न 18
के साथ दोनों \omterm{def:g12:seq:sequence}{अनुक्रम} संलग्न हैं तथा एक सीमा $M(1, 2)$ साझा करते हैं।

\textbf{20.} तीसरी पुनरावृत्ति $a_3 \approx b_3 \approx 1.456791$ देती है: तीन ही घुमावों में $M(1, 2) \approx
1.456791$ (अंतर
मोटे तौर पर \emph{वर्ग} हो जाता है, हीरोन की तरह)। इस समस्या के \omterm{def:g11:seq:sequence}{अनुक्रमों} की
गति-क्रम में सूची, सबसे धीमे से सबसे तेज़ तक: हरात्मक योग (हिमनद-सा अपसरण),
\omterm{def:g12:seq:arith-geom}{गुणोत्तर} योग (हर क़दम पर त्रुटि आधी), हीरोन और समांतर--गुणोत्तर \omterm{def:g11:stat:mean}{माध्य} (हर
क़दम पर त्रुटि का वर्ग) --- और उसी \omterm{def:g11:stat:mean}{माध्य} की अलौकिक गति ने गाउस को बता दिया था
कि उन्होंने विश्लेषण की नई खान खोद ली है।
\end{solution}
